% ═══════════════════════════════════════════════════════════════
% ADICIONAR AO PREÂMBULO DO DOCUMENTO PRINCIPAL:
%
% \usepackage{listings}
% \usepackage{xcolor}
% \lstset{
%   basicstyle=\ttfamily\small,
%   keywordstyle=\color{blue}\bfseries,
%   commentstyle=\color{gray}\itshape,
%   stringstyle=\color{orange},
%   showstringspaces=false,
%   breaklines=true,
%   frame=single,
%   language=Python,
%   numbers=left,
%   numberstyle=\tiny\color{gray},
%   xleftmargin=1.5em,
%   framexleftmargin=1.5em,
% }
% ═══════════════════════════════════════════════════════════════

\section{Implementação em Python}
\label{sec:implementacao}

Esta secção descreve a implementação prática de cada método em Python,
assumindo que se dispõe de dois ficheiros de malha separados ---
endocárdio (\texttt{endo}) e epicárdio (\texttt{epi}) --- carregados
com a biblioteca \texttt{trimesh}, e que o domínio volumétrico é
representado por máscaras binárias NumPy para os métodos baseados em
voxels.

\subsection*{Configuração inicial comum}

O bloco seguinte é partilhado por todos os métodos.
As malhas são carregadas e as máscaras binárias são obtidas
a partir das mesmas por voxelização.

\begin{lstlisting}
import numpy as np
import trimesh
from scipy.spatial import cKDTree
from scipy.ndimage import distance_transform_edt, label, maximum_filter

# ── Carregar malhas ──────────────────────────────────────────────
endo = trimesh.load("endo.stl")   # superficie interior
epi  = trimesh.load("epi.stl")    # superficie exterior

P = np.array(endo.vertices)   # (N, 3) pontos do endocardio
Q = np.array(epi.vertices)    # (M, 3) pontos do epicardio
N_endo = np.array(endo.vertex_normals)  # normais por vertice

# ── Voxelizacao (para metodos volumetricos) ──────────────────────
# pitch = resolucao em mm por voxel
pitch = 1.0
vol_endo = endo.voxelized(pitch=pitch).matrix   # mascara binaria
vol_epi  = epi.voxelized(pitch=pitch).matrix

# mascara do miocardio = entre as duas superficies
# (assumindo que epi envolve endo)
vol_myo = vol_epi.astype(bool) & ~vol_endo.astype(bool)
\end{lstlisting}

% ───────────────────────────────────────────────────────────────
\subsection{Método 1 --- Vizinho Mais Próximo (KD-tree)}
\label{impl:kdtree}
% ───────────────────────────────────────────────────────────────

\paragraph{Como funciona.}
Constrói-se uma KD-tree com os vértices do epicárdio.
Para cada vértice do endocárdio, a \emph{query} à árvore devolve
a distância ao ponto mais próximo do epicárdio em tempo
$O(\log M)$, onde $M$ é o número de vértices do epi.

\begin{lstlisting}
from scipy.spatial import cKDTree

tree_epi = cKDTree(Q)

# Para cada ponto do endo, distancia ao vizinho mais proximo no epi
t_kdtree, _ = tree_epi.query(P, workers=-1)  # workers=-1 usa todos os cores

# t_kdtree : array (N,) com a espessura estimada por vertice
print(f"Espessura media: {t_kdtree.mean():.2f} mm")
print(f"Espessura min/max: {t_kdtree.min():.2f} / {t_kdtree.max():.2f} mm")
\end{lstlisting}

\paragraph{Notas de implementação.}
O parâmetro \texttt{workers=-1} paraleliza automaticamente a query em
todos os núcleos disponíveis. Para malhas com mais de 100\,000 vértices,
a diferença de tempo é significativa. O resultado é válido como
\emph{baseline} mas subestima a espessura no ápex e em zonas curvas.

% ───────────────────────────────────────────────────────────────
\subsection{Método 2 --- KD-tree Simétrico}
\label{impl:symkd}
% ───────────────────────────────────────────────────────────────

\paragraph{Como funciona.}
Executa-se a query KD-tree nos dois sentidos: endo$\to$epi
e epi$\to$endo. A diferença entre as médias ($\Delta$) é um
indicador de qualidade --- se for grande, o método está a
introduzir enviesamento geométrico.

\begin{lstlisting}
tree_endo = cKDTree(P)
tree_epi  = cKDTree(Q)

# Endo -> Epi
t_endo2epi, _ = tree_epi.query(P, workers=-1)

# Epi -> Endo
t_epi2endo, _ = tree_endo.query(Q, workers=-1)

delta = abs(t_endo2epi.mean() - t_epi2endo.mean())
print(f"Delta simetria: {delta:.3f} mm")

if delta < 0.5:
    print("OK: os dois sentidos concordam.")
else:
    print("AVISO: enviesamento geometrico detectado.")

# O valor de espessura usado e a media dos dois sentidos
t_sym = 0.5 * (t_endo2epi + t_epi2endo[
    tree_endo.query(P, workers=-1)[1]
])
\end{lstlisting}

\paragraph{Notas de implementação.}
Este método serve principalmente para diagnóstico.
Se $\Delta > 1$\,mm numa malha cardíaca típica, a qualidade
da malha deve ser verificada antes de prosseguir.

% ───────────────────────────────────────────────────────────────
\subsection{Método 3 --- Raios Normais}
\label{impl:rays}
% ───────────────────────────────────────────────────────────────

\paragraph{Como funciona.}
Para cada vértice do endocárdio, lança-se um raio na direção
da normal à superfície. A intersecção do raio com o epicárdio
dá a espessura local. O \texttt{trimesh} usa uma estrutura BVH
(\emph{Bounding Volume Hierarchy}) internamente, tornando o
processo eficiente mesmo para malhas densas.

\begin{lstlisting}
import trimesh

# Criador de intersetor BVH sobre o epicardio
intersector = trimesh.ray.ray_triangle.RayMeshIntersector(epi)

origins    = P                # (N, 3) pontos de partida
directions = N_endo           # (N, 3) normais do endocardio

# Lancar todos os raios de uma vez
locations, index_ray, _ = intersector.intersects_location(
    ray_origins=origins,
    ray_directions=directions,
    multiple_hits=False       # apenas o primeiro impacto
)

# Calcular espessura: distancia do ponto de partida ao impacto
t_rays = np.full(len(P), np.nan)   # NaN onde o raio falhou

for loc, idx in zip(locations, index_ray):
    t_rays[idx] = np.linalg.norm(loc - P[idx])

# Tratar raios falhados: preencher com interpolacao dos vizinhos
failed = np.isnan(t_rays)
print(f"Raios falhados: {failed.sum()} / {len(P)} "
      f"({100*failed.mean():.1f}%)")

# Interpolacao simples para raios falhados
from scipy.spatial import cKDTree
tree_ok = cKDTree(P[~failed])
_, nn_idx = tree_ok.query(P[failed])
t_rays[failed] = t_rays[~failed][nn_idx]
\end{lstlisting}

\paragraph{Notas de implementação.}
Uma taxa de raios falhados acima de 5\% indica normais invertidas
ou gaps na malha do epicárdio. Nesse caso, deve verificar-se a
orientação das normais com \texttt{trimesh.repair.fix\_normals(epi)}.
Para malhas muito ruidosas, considere o Método~8 (SDF cone) em vez
deste.

% ───────────────────────────────────────────────────────────────
\subsection{Método 4 --- EDT Soma de Fronteiras}
\label{impl:edt}
% ───────────────────────────────────────────────────────────────

\paragraph{Como funciona.}
A \texttt{scipy} calcula a Transformada de Distância Euclidiana
(EDT) de forma eficiente em 3D. Calculam-se duas EDT --- uma
a partir da fronteira do endocárdio e outra da fronteira do
epicárdio --- e somam-se dentro do miocárdio.

\begin{lstlisting}
from scipy.ndimage import distance_transform_edt

# Fronteiras binarias (True = fronteira, False = interior)
# A EDT calcula a distancia ao ponto True mais proximo
# Invertemos: passamos a mascara do interior de cada regiao
d_endo = distance_transform_edt(~vol_endo) * pitch  # mm
d_epi  = distance_transform_edt(~vol_epi)  * pitch  # mm

# Espessura em cada voxel do miocardio
t_edt_vol = (d_endo + d_epi) * vol_myo.astype(float)

# Espessura media no miocardio
t_edt_mean = t_edt_vol[vol_myo].mean()
print(f"Espessura media (EDT): {t_edt_mean:.2f} mm")

# Projectar de volta para os vertices do endocardio
# (convertendo coordenadas do mundo para indices de voxel)
def world_to_voxel(pts, origin, pitch):
    return np.round((pts - origin) / pitch).astype(int)

origin = endo.bounds[0]  # canto inferior esquerdo da bounding box
vox_idx = world_to_voxel(P, origin, pitch)
vox_idx = np.clip(vox_idx, 0,
                  np.array(t_edt_vol.shape) - 1)

t_edt = t_edt_vol[vox_idx[:,0], vox_idx[:,1], vox_idx[:,2]]
\end{lstlisting}

\paragraph{Notas de implementação.}
O factor \texttt{pitch} converte de voxels para milímetros.
A resolução da imagem ($\leq 1$\,mm/voxel) é o factor limitante
da precisão. Para imagens com resolução anisotrópica,
\texttt{distance\_transform\_edt} aceita o argumento
\texttt{sampling=(dz, dy, dx)}.

% ───────────────────────────────────────────────────────────────
\subsection{Método 5 --- EDT Medial (Eixo Medial)}
\label{impl:medial}
% ───────────────────────────────────────────────────────────────

\paragraph{Como funciona.}
O eixo medial corresponde aos máximos locais da EDT dentro do
miocárdio. Esses pontos de ``crista'' têm pelo menos dois vizinhos
de fronteira à mesma distância. O valor da EDT nesses pontos é
aproximadamente metade da espessura local.

\begin{lstlisting}
from scipy.ndimage import maximum_filter

# EDT dentro do miocardio (distancia a qualquer fronteira)
d_any = distance_transform_edt(vol_myo) * pitch

# Maximos locais numa vizinhanca 3x3x3
local_max = maximum_filter(d_any, size=3)
is_ridge  = (d_any == local_max) & vol_myo & (d_any > 0)

# Espessura no eixo medial
t_medial_pts = 2 * d_any[is_ridge]
print(f"Espessura media no eixo medial: {t_medial_pts.mean():.2f} mm")
print(f"Numero de pontos de crista: {is_ridge.sum()}")

# Visualizacao: coordenadas 3D dos pontos de crista
origin = endo.bounds[0]
ridge_coords = np.argwhere(is_ridge) * pitch + origin
\end{lstlisting}

\paragraph{Notas de implementação.}
Este método produz poucos pontos e é mais adequado para
verificação rápida do que para um mapa denso. Para um mapa mais
denso, usar o Método~4 ou o Método~3.

% ───────────────────────────────────────────────────────────────
\subsection{Método 6 --- Campo de Laplace}
\label{impl:laplace}
% ───────────────────────────────────────────────────────────────

\paragraph{Como funciona.}
Resolve-se $\nabla^2\psi = 0$ dentro do miocárdio com condições
de fronteira $\psi=0$ no endocárdio e $\psi=1$ no epicárdio.
O sistema linear esparso resultante é resolvido com um solver
iterativo (\texttt{scipy.sparse.linalg}). A espessura local é
$t \approx 1/|\nabla\psi|$ avaliada em $\psi \approx 0.5$.

\begin{lstlisting}
import scipy.sparse as sp
import scipy.sparse.linalg as spla

# Construir o sistema linear para a equacao de Laplace
# nos voxels do miocardio (stencil de 6 vizinhos em 3D)

shape = vol_myo.shape
idx_myo = np.argwhere(vol_myo)      # (K, 3) indices dos voxels do myo
n = len(idx_myo)

# Mapa de indice global para indice no sistema linear
global_to_local = -np.ones(shape, dtype=int)
for local_i, gidx in enumerate(idx_myo):
    global_to_local[tuple(gidx)] = local_i

rows, cols, vals = [], [], []
b = np.zeros(n)

neighbors = [(1,0,0),(-1,0,0),(0,1,0),(0,-1,0),(0,0,1),(0,0,-1)]

for local_i, (iz, iy, ix) in enumerate(idx_myo):
    rows.append(local_i); cols.append(local_i); vals.append(6.0)

    for dz, dy, dx in neighbors:
        nz, ny, nx = iz+dz, iy+dy, ix+dx

        if not (0 <= nz < shape[0] and
                0 <= ny < shape[1] and
                0 <= nx < shape[2]):
            continue

        if vol_endo[nz, ny, nx]:
            b[local_i] += 0.0   # fronteira endo: psi = 0
        elif not vol_myo[nz, ny, nx]:
            # fronteira epi (fora do myo): psi = 1
            b[local_i] += 1.0
        else:
            j = global_to_local[nz, ny, nx]
            rows.append(local_i); cols.append(j); vals.append(-1.0)

A = sp.csr_matrix((vals, (rows, cols)), shape=(n, n))

# Resolver o sistema (CG = Gradiente Conjugado, rapido para Laplace)
psi_vec, info = spla.cg(A, b, tol=1e-5, maxiter=5000)
print(f"Solver convergiu: {info == 0}")

# Reconstruir volume psi
psi_vol = np.zeros(shape)
for local_i, gidx in enumerate(idx_myo):
    psi_vol[tuple(gidx)] = psi_vec[local_i]

# Gradiente e espessura
dpsi_dz, dpsi_dy, dpsi_dx = np.gradient(psi_vol, pitch)
grad_mag = np.sqrt(dpsi_dz**2 + dpsi_dy**2 + dpsi_dx**2)

# Evitar divisao por zero
grad_mag[grad_mag < 1e-6] = np.nan
t_laplace_vol = 1.0 / grad_mag
t_laplace_vol[~vol_myo] = np.nan

# Projectar para vertices do endocardio
origin = endo.bounds[0]
vox_idx = world_to_voxel(P, origin, pitch)
vox_idx = np.clip(vox_idx, 0, np.array(shape) - 1)
t_laplace = t_laplace_vol[vox_idx[:,0],
                           vox_idx[:,1],
                           vox_idx[:,2]]
\end{lstlisting}

\paragraph{Notas de implementação.}
Para volumes grandes ($>$ 200$^3$ voxels), o solver CG pode ser
lento. Alternativas mais rápidas incluem o solver
\texttt{pyamg} (multigrid algébrico) ou reduzir o
\texttt{pitch} para 2\,mm. O método \texttt{pyezzi}
(Secção~\ref{impl:yezzi}) é uma alternativa mais eficiente
com qualidade equivalente.

% ───────────────────────────────────────────────────────────────
\subsection{Método 7 --- Geodésico / Dijkstra}
\label{impl:geodesic}
% ───────────────────────────────────────────────────────────────

\paragraph{Como funciona.}
Cada voxel do miocárdio é um nó de um grafo. As arestas ligam
voxels vizinhos (6 ou 26 vizinhança) com pesos iguais à distância
euclidiana entre centros. O algoritmo de Dijkstra (implementado
como \emph{Bellman-Ford} esparso no \texttt{scipy}) calcula
a distância mínima de todos os voxels às fronteiras.

\begin{lstlisting}
from scipy.sparse.csgraph import dijkstra
import scipy.sparse as sp

# Construir grafo esparso dos voxels do miocardio
idx_myo = np.argwhere(vol_myo)
n = len(idx_myo)
global_to_local = -np.ones(shape, dtype=int)
for li, gidx in enumerate(idx_myo):
    global_to_local[tuple(gidx)] = li

row_g, col_g, w_g = [], [], []
neighbors_26 = [(dz, dy, dx)
                for dz in [-1,0,1]
                for dy in [-1,0,1]
                for dx in [-1,0,1]
                if not (dz==dy==dx==0)]

for li, (iz, iy, ix) in enumerate(idx_myo):
    for dz, dy, dx in neighbors_26:
        nz, ny, nx = iz+dz, iy+dy, ix+dx
        if (0 <= nz < shape[0] and
            0 <= ny < shape[1] and
            0 <= nx < shape[2] and
            vol_myo[nz, ny, nx]):
            lj = global_to_local[nz, ny, nx]
            dist_vox = np.sqrt(dz**2 + dy**2 + dx**2) * pitch
            row_g.append(li); col_g.append(lj); w_g.append(dist_vox)

G = sp.csr_matrix((w_g, (row_g, col_g)), shape=(n, n))

# Nos de fronteira endo e epi
endo_nodes = [global_to_local[tuple(i)]
              for i in idx_myo
              if any(vol_endo[tuple(i + np.array(nb))]
                     for nb in [(1,0,0),(-1,0,0),(0,1,0),
                                (0,-1,0),(0,0,1),(0,0,-1)]
                     if 0 <= i[0]+nb[0] < shape[0])]

epi_nodes  = [global_to_local[tuple(i)]
              for i in idx_myo
              if not any(vol_myo[tuple(i + np.array(nb))]
                         for nb in [(1,0,0),(-1,0,0),(0,1,0),
                                    (0,-1,0),(0,0,1),(0,0,-1)]
                         if 0 <= i[0]+nb[0] < shape[0])]

# Dijkstra a partir das fronteiras
d_to_endo = dijkstra(G, indices=endo_nodes).min(axis=0)
d_to_epi  = dijkstra(G, indices=epi_nodes).min(axis=0)

t_geo_vec = d_to_endo + d_to_epi

# Reconstruir volume
t_geo_vol = np.zeros(shape)
for li, gidx in enumerate(idx_myo):
    t_geo_vol[tuple(gidx)] = t_geo_vec[li]
\end{lstlisting}

\paragraph{Notas de implementação.}
A vizinhança de 26 nós (vs.\ 6) reduz o artefacto de
discretização na diagonal, mas aumenta o custo computacional.
Para volumes grandes, o Dijkstra esparso do \texttt{scipy} pode
ser lento; considere usar o Método~4 (EDT) como aproximação.

% ───────────────────────────────────────────────────────────────
\subsection{Método 8 --- SDF / Cone de Raios}
\label{impl:sdf}
% ───────────────────────────────────────────────────────────────

\paragraph{Como funciona.}
Em vez de um único raio por vértice (Método~3), lançam-se $K$
raios dentro de um cone de abertura $\alpha$ à volta da normal.
A mediana das distâncias de impacto é mais robusta a normais
ruidosas e a triângulos mal orientados.

\begin{lstlisting}
def cone_rays(normal, K=7, alpha_deg=30.0):
    """Gera K direccoes dentro de um cone de abertura alpha."""
    alpha = np.radians(alpha_deg)
    # Base ortonormal local
    perp = np.array([1, 0, 0]) if abs(normal[0]) < 0.9 \
           else np.array([0, 1, 0])
    u = np.cross(normal, perp); u /= np.linalg.norm(u)
    v = np.cross(normal, u)

    phi = np.linspace(0, 2*np.pi, K, endpoint=False)
    dirs = (normal * np.cos(alpha) +
            np.outer(np.cos(phi), u) * np.sin(alpha) +
            np.outer(np.sin(phi), v) * np.sin(alpha))
    return dirs / np.linalg.norm(dirs, axis=1, keepdims=True)

intersector = trimesh.ray.ray_triangle.RayMeshIntersector(epi)
K = 7

t_sdf = np.full(len(P), np.nan)

for i, (pt, nrm) in enumerate(zip(P, N_endo)):
    dirs = cone_rays(nrm, K=K, alpha_deg=30)
    origins_cone = np.tile(pt, (K, 1))

    locs, idx_ray, _ = intersector.intersects_location(
        ray_origins=origins_cone,
        ray_directions=dirs,
        multiple_hits=False
    )

    if len(locs) == 0:
        continue

    dists = np.linalg.norm(locs - pt, axis=1)
    t_sdf[i] = np.median(dists)

# Tratar falhas com interpolacao
failed = np.isnan(t_sdf)
if failed.any():
    tree_ok = cKDTree(P[~failed])
    _, nn = tree_ok.query(P[failed])
    t_sdf[failed] = t_sdf[~failed][nn]

print(f"SDF media: {np.nanmean(t_sdf):.2f} mm")
\end{lstlisting}

\paragraph{Notas de implementação.}
$K=7$ e $\alpha=30°$ são os valores recomendados pela biblioteca
CGAL para geometrias semelhantes ao VE. Aumentar $K$ melhora
a robustez mas aumenta o tempo linear em $K$. Para malhas
limpas, $K=5$ é suficiente.

% ───────────────────────────────────────────────────────────────
\subsection{Método 9 --- Correspondência Regularizada}
\label{impl:reg}
% ───────────────────────────────────────────────────────────────

\paragraph{Como funciona.}
Começa com uma estimativa inicial baseada na projeção na normal
(semelhante ao Método~3 mas usando o KD-tree como aproximação).
Depois aplica um filtro de suavização iterativo que combina o valor
local com a média dos vizinhos na malha, produzindo um campo de
espessura suave e sem saltos.

\begin{lstlisting}
# ── Passo 1: estimativa inicial ──────────────────────────────────
tree_epi = cKDTree(Q)
_, nn_idx = tree_epi.query(P)
q_nearest = Q[nn_idx]                           # (N, 3)

proj   = np.sum((q_nearest - P) * N_endo, axis=1)  # projecao na normal
dist   = np.linalg.norm(q_nearest - P, axis=1)
t0     = np.maximum(proj, dist)                 # estimativa inicial

# ── Passo 2: grafo de adjacencia da malha do endocardio ──────────
adjacency = {i: set() for i in range(len(P))}
for face in endo.faces:
    a, b, c = face
    adjacency[a].update([b, c])
    adjacency[b].update([a, c])
    adjacency[c].update([a, b])

# ── Passo 3: suavizacao iterativa ────────────────────────────────
alpha  = 0.35
n_iter = 10
t_reg  = t0.copy()

for iteration in range(n_iter):
    t_new = np.zeros_like(t_reg)
    for i in range(len(P)):
        nbrs = list(adjacency[i])
        if nbrs:
            t_nbr_mean = t_reg[nbrs].mean()
        else:
            t_nbr_mean = t_reg[i]
        t_new[i] = (1 - alpha) * t0[i] + alpha * t_nbr_mean
    t_reg = t_new
    if (iteration + 1) % 5 == 0:
        print(f"  iter {iteration+1}: mean={t_reg.mean():.3f} mm")

print(f"Espessura regularizada media: {t_reg.mean():.2f} mm")
\end{lstlisting}

\paragraph{Notas de implementação.}
O parâmetro $\alpha=0.35$ é o valor recomendado para malhas
cardíacas. Valores de $\alpha$ superiores a 0.6 podem
sobre-suavizar e esconder zonas de enfarte. O número de iterações
($n=10$) pode ser aumentado até estabilização (convergência
tipicamente antes de 20 iterações).

% ───────────────────────────────────────────────────────────────
\subsection{Método 10 --- Yezzi--Prince (PDE Euleriana)}
\label{impl:yezzi}
% ───────────────────────────────────────────────────────────────

\paragraph{Como funciona.}
Yezzi e Prince~\cite{yezzi2003} propuseram uma formulação PDE
alternativa à equação de Laplace. Em vez de resolver uma equação,
resolve-se um par de PDEs guiadas pelo campo vetorial transmural:
\begin{align}
  \nabla\psi \cdot \nabla u &= -1, \quad u\big|_{\text{endo}} = 0 \\
  \nabla\psi \cdot \nabla v &= -1, \quad v\big|_{\text{epi}} = 0
\end{align}
A espessura é $t = u + v$, sem necessidade de construir
trajetórias explícitas. O método é mais estável numericamente
que o Laplace puro, em especial no ápex e nas zonas de parede
fina. A implementação Python está disponível no pacote
\texttt{pyezzi}.

\begin{lstlisting}
# Instalar: pip install pyezzi

from pyezzi import compute_thickness_cardiac

# vol_endo e vol_epi sao mascaras binarias NumPy (bool)
# endo = interior (cavidade), epi = exterior (incluindo o myo)
t_yezzi_vol = compute_thickness_cardiac(
    endo=vol_endo,
    epi=vol_epi
) * pitch  # converter de voxels para mm

# Espessura media
print(f"Espessura media (Yezzi-Prince): "
      f"{t_yezzi_vol[vol_myo].mean():.2f} mm")

# Projectar para vertices do endocardio
origin  = endo.bounds[0]
vox_idx = world_to_voxel(P, origin, pitch)
vox_idx = np.clip(vox_idx, 0, np.array(t_yezzi_vol.shape) - 1)
t_yezzi = t_yezzi_vol[vox_idx[:,0],
                       vox_idx[:,1],
                       vox_idx[:,2]]
\end{lstlisting}

\paragraph{Notas de implementação.}
O \texttt{pyezzi} aceita máscaras binárias anisotrópicas através
do argumento \texttt{sampling=(dz, dy, dx)} (em mm), o que é
útil para dados de MRI com resolução não uniforme nos três eixos.
É o método recomendado quando se pretende qualidade equivalente
ao Laplace com menor custo computacional.

% ───────────────────────────────────────────────────────────────
\subsection{Mapa de Cor 3D e Mapa AHA-17}
\label{impl:colormap}
% ───────────────────────────────────────────────────────────────

\paragraph{Como funciona.}
Após calcular o array de espessura por vértice com qualquer
dos métodos anteriores, a visualização 3D é obtida mapeando
os valores para uma escala de cor contínua. O mapa AHA-17
é obtido agregando os valores por segmento.

\begin{lstlisting}
import matplotlib.pyplot as plt
import matplotlib.colors as mcolors

# ── Mapa de cor 3D ───────────────────────────────────────────────
# Usar t_rays, t_sdf, t_yezzi, etc.
t = t_rays.copy()                    # escolher o metodo
t = np.nan_to_num(t, nan=np.nanmedian(t))

# Normalizar para [0, 1]
t_min, t_max = np.percentile(t, [2, 98])   # percentis para robustez
t_norm = np.clip((t - t_min) / (t_max - t_min), 0, 1)

# Aplicar colormap (ex: 'RdYlGn_r': vermelho=fino, verde=normal)
cmap   = plt.cm.RdYlGn_r
colors = (cmap(t_norm)[:, :3] * 255).astype(np.uint8)

# Atribuir cores aos vertices do endocardio
endo.visual = trimesh.visual.ColorVisuals(
    mesh=endo,
    vertex_colors=np.hstack([colors,
                              255*np.ones((len(P),1), dtype=np.uint8)])
)
endo.export("endo_thickness_colormap.ply")   # abrir no MeshLab/Blender

# ── Mapa AHA-17 (bull's-eye) ─────────────────────────────────────
# Requer: aha_segment[i] = segmento AHA (1..17) de cada vertice i
# (calculado separadamente com base no eixo longo + pontos RV)

def compute_aha_map(t, aha_segment):
    """Calcula a espessura media por segmento AHA."""
    aha_means = {}
    for seg in range(1, 18):
        mask = aha_segment == seg
        if mask.sum() > 0:
            aha_means[seg] = t[mask].mean()
        else:
            aha_means[seg] = np.nan
    return aha_means

# aha_segment = array (N,) com o segmento de cada vertice
# aha_map = compute_aha_map(t, aha_segment)
\end{lstlisting}

% ───────────────────────────────────────────────────────────────
\subsection{Comparação rápida dos resultados}
\label{impl:compare}
% ───────────────────────────────────────────────────────────────

O bloco seguinte agrega todos os métodos numa tabela de comparação,
útil para incluir na tese como validação cruzada.

\begin{lstlisting}
methods = {
    "KD-tree":          t_kdtree,
    "KD-tree simetrico":t_sym,
    "Raios normais":    t_rays,
    "SDF cone":         t_sdf,
    "Corr. regularizada": t_reg,
    "Yezzi-Prince":     t_yezzi,
}

print(f"{'Metodo':<25} {'Media':>8} {'Std':>8} "
      f"{'Min':>8} {'Max':>8}")
print("-" * 60)
for name, t_arr in methods.items():
    valid = t_arr[~np.isnan(t_arr)]
    print(f"{name:<25} {valid.mean():>8.2f} {valid.std():>8.2f} "
          f"{valid.min():>8.2f} {valid.max():>8.2f}")
\end{lstlisting}

% ───────────────────────────────────────────────────────────────
\subsection*{Referência}
% ───────────────────────────────────────────────────────────────

\begin{enumerate}[label={[\arabic*]}, nosep]
  \item Yezzi, A.J., Prince, J.L. (2003). An Eulerian PDE approach
        for computing tissue thickness. \emph{IEEE Trans.\ Med.\
        Imaging}, 22(10):1332–1339.
        \label{yezzi2003}
  \item Jones, S.E., Buchbinder, B.R., Aharon, I. (2000).
        Three-dimensional mapping of cortical thickness using
        Laplace's equation. \emph{Human Brain Mapping},
        11(1):12–32.
  \item Osher, S., Sethian, J.A. (1988). Shape Diameter Function.
        \emph{CGAL Computational Geometry Algorithms Library}.
\end{enumerate}
\end{document}